% TOP_PROBLEM: {"id": 1501, "title": "Infinitude of Amicable Pairs", "category_id": 1, "category": {"id": 1, "name": "number_theory", "display_name": "Number Theory", "description": "Properties of integers, prime numbers, Diophantine equations.", "slug": "number-theory", "order_index": 1, "created_at": "2026-07-31T15:26:25.670Z"}, "status": "open"}
% FIRST_POSTED: null
% ENTRY_KIND: statement_only
% Imported from: batch08.tex; UnsolvedMath ID 1501
\documentclass[11pt]{article}
\usepackage[margin=25mm]{geometry}
\usepackage{fontspec}
\setmainfont{Latin Modern Roman}
\usepackage{amsmath,amssymb,mathtools}
\usepackage{unicode-math}
\setmathfont{Latin Modern Math}
\usepackage{xurl,hyperref}
\hypersetup{colorlinks=true,urlcolor=blue}
\setcounter{secnumdepth}{0}
\setlength{\parindent}{0pt}
\setlength{\parskip}{5pt}
\setlength{\emergencystretch}{3em}
\newcommand{\sref}[2]{\hyperlink{source:#1:#2}{[S#2]}}
\newcommand{\cataloguescope}[1]{\paragraph{Recorded target.}#1}
\begin{document}
% UnsolvedMath ID 1501; source catalogue code NUM-007
\setcounter{section}{1500}
\section{Infinitude of Amicable Pairs}\label{1501}
{\small\sffamily UnsolvedMath ID 1501 · Number Theory}
\subsection{Definitions and mathematical statement}

\paragraph{Shared notation.}$\mathbb N=\{1,2,\ldots\}$, and $\mathbb Z,\mathbb Q,\mathbb R,\mathbb C$ denote the integers, rationals, reals and complex numbers. Any other conventions are specified locally.


Let \(\sigma(n)=\sum_{d\mid n,\,d>0}d\), and let
\(s(n)=\sigma(n)-n\) be the sum of the proper positive divisors.
An amicable pair consists of distinct positive integers \(a,b\)
with \(s(a)=b\) and \(s(b)=a\). Represent each unordered pair
once by choosing \(a<b\). No restriction on the parities is imposed.
\paragraph{Statement recorded in the supplied source.}
Is the set
\[
 \{(a,b)\in\mathbb N^2:a<b,\ \sigma(a)-a=b,\ \sigma(b)-b=a\}
\]
infinite? \sref{1501}{2}
\subsection{Short English statement}
Are there infinitely many distinct pairs in which each integer is the sum of the other's proper positive divisors?
\subsection{Sources}
\begin{description}
\item[\hypertarget{source:1501:1}{S1}] ULAM.AI and the UnsolvedMath Contributors, \emph{UnsolvedMath: A Curated Collection of Open Mathematics Problems}, record 1501 (NUM-007). CC BY 4.0. \url{https://huggingface.co/datasets/ulamai/UnsolvedMath}.
\item[\hypertarget{source:1501:2}{S2}] Germano D'Abramo, \emph{On Amicable Numbers With Different Parity}, arXiv:math/0501402v4 (2007), Section 1, definition of amicable numbers and the general infinitude question. \url{https://arxiv.org/abs/math/0501402}.
\end{description}
\label{end:1501}
\end{document}
